\section{Morphisms of finite type}
\label{section-finite-type}

\noindent
Recall that a ring map $R \to A$ is said to be of finite type if
$A$ is isomorphic to a quotient of $R[x_1, \ldots, x_n]$ as an $R$-algebra, see
Algebra, Definition \ref{algebra-definition-finite-type}.

\begin{definition}
\label{definition-finite-type}
Let $f : X \to S$ be a morphism of schemes.
\begin{enumerate}
\item We say that $f$ is of {\it finite type at $x \in X$} if
there exists an affine open neighbourhood $\Spec(A) = U \subset X$
of $x$ and an affine open $\Spec(R) = V \subset S$
with $f(U) \subset V$ such that the induced ring map
$R \to A$ is of finite type.
\item We say that $f$ is {\it locally of finite type} if it is
of finite type at every point of $X$.
\item We say that $f$ is of {\it finite type} if it is locally of
finite type and quasi-compact.
\end{enumerate}
\end{definition}

\begin{lemma}
\label{lemma-locally-finite-type-characterize}
Let $f : X \to S$ be a morphism of schemes.
The following are equivalent
\begin{enumerate}
\item The morphism $f$ is locally of finite type.
\item For all affine opens $U \subset X$, $V \subset S$
with $f(U) \subset V$ the ring map
$\mathcal{O}_S(V) \to \mathcal{O}_X(U)$ is of finite type.
\item There exist an open covering $S = \bigcup_{j \in J} V_j$
and open coverings $f^{-1}(V_j) = \bigcup_{i \in I_j} U_i$ such
that each of the morphisms $U_i \to V_j$, $j\in J, i\in I_j$
is locally of finite type.
\item There exist an affine open covering $S = \bigcup_{j \in J} V_j$
and affine open coverings $f^{-1}(V_j) = \bigcup_{i \in I_j} U_i$ such
that the ring map $\mathcal{O}_S(V_j) \to \mathcal{O}_X(U_i)$ is
of finite type, for all $j\in J, i\in I_j$.
\end{enumerate}
Moreover, if $f$ is locally of finite type then for
any open subschemes $U \subset X$, $V \subset S$ with $f(U) \subset V$
the restriction $f|_U : U \to V$ is locally of finite type.
\end{lemma}

\begin{proof}
This follows from Lemma \ref{lemma-locally-P} if we show that
the property ``$R \to A$ is of finite type'' is local.
We check conditions (a), (b) and (c) of Definition
\ref{definition-property-local}.
By Algebra, Lemma \ref{algebra-lemma-base-change-finiteness}
being of finite type is stable under base change and hence
we conclude (a) holds. By Algebra, Lemma
\ref{algebra-lemma-compose-finite-type} being of finite type
is stable under composition and trivially for any ring
$R$ the ring map $R \to R_f$ is of finite type.
We conclude (b) holds. Finally, property (c) is true
according to Algebra, Lemma \ref{algebra-lemma-cover-upstairs}.
\end{proof}

\begin{lemma}
\label{lemma-composition-finite-type}
The composition of two morphisms which are locally of finite type is
locally of finite type. The same is true for morphisms of finite type.
\end{lemma}

\begin{proof}
In the proof of Lemma \ref{lemma-locally-finite-type-characterize}
we saw that being of finite type is a local property of ring maps.
Hence the first statement of the lemma follows from
Lemma \ref{lemma-composition-type-P} combined
with the fact that being of finite type is a property of ring maps that is
stable under composition, see
Algebra, Lemma \ref{algebra-lemma-compose-finite-type}.
By the above and the fact that compositions of
quasi-compact morphisms are quasi-compact, see
Schemes, Lemma \ref{schemes-lemma-composition-quasi-compact}
we see that the composition of morphisms of finite type is
of finite type.
\end{proof}

\begin{lemma}
\label{lemma-base-change-finite-type}
The base change of a morphism which is locally of finite type
is locally of finite type. The same is true for morphisms of
finite type.
\end{lemma}

\begin{proof}
In the proof of Lemma \ref{lemma-locally-finite-type-characterize}
we saw that being of finite type is a local property of ring maps.
Hence the first statement of the lemma follows from
Lemma \ref{lemma-base-change-type-P} combined
with the fact that being of finite type is a property of ring maps that is
stable under base change, see
Algebra, Lemma \ref{algebra-lemma-base-change-finiteness}.
By the above and the fact that a base change of a
quasi-compact morphism is quasi-compact, see
Schemes, Lemma \ref{schemes-lemma-quasi-compact-preserved-base-change}
we see that the base change of a morphism of finite type is
a morphism of finite type.
\end{proof}

\begin{lemma}
\label{lemma-immersion-locally-finite-type}
A closed immersion is of finite type.
An immersion is locally of finite type.
\end{lemma}

\begin{proof}
This is true because an open immersion is a local isomorphism,
and a closed immersion is obviously of finite type.
\end{proof}

\begin{lemma}
\label{lemma-finite-type-noetherian}
Let $f : X \to S$ be a morphism.
If $S$ is (locally) Noetherian and $f$ (locally) of finite type
then $X$ is (locally) Noetherian.
\end{lemma}

\begin{proof}
This follows immediately from the fact that a ring
of finite type over a Noetherian ring is Noetherian,
see Algebra, Lemma \ref{algebra-lemma-Noetherian-permanence}.
(Also: use the fact that the source of a quasi-compact morphism
with quasi-compact target is quasi-compact.)
\end{proof}

\begin{lemma}
\label{lemma-finite-type-Noetherian-quasi-separated}
Let $f : X \to S$ be locally of finite type with $S$ locally Noetherian.
Then $f$ is quasi-separated.
\end{lemma}

\begin{proof}
In fact, it is true that $X$ is quasi-separated, see
Properties, Lemma \ref{properties-lemma-locally-Noetherian-quasi-separated}
and Lemma \ref{lemma-finite-type-noetherian} above.
Then apply Schemes, Lemma \ref{schemes-lemma-compose-after-separated}
to conclude that $f$ is quasi-separated.
\end{proof}

\begin{lemma}
\label{lemma-permanence-finite-type}
Let $X \to Y$ be a morphism of schemes over a base scheme $S$.
If $X$ is locally of finite type over $S$, then $X \to Y$
is locally of finite type.
\end{lemma}

\begin{proof}
Via Lemma \ref{lemma-locally-finite-type-characterize} this translates
into the following algebra
fact: Given ring maps $A \to B \to C$ such that $A \to C$ is
of finite type, then $B \to C$ is of finite type.
(See
Algebra, Lemma \ref{algebra-lemma-compose-finite-type}).
\end{proof}







